\label{sec:limitations}

We restrict our study to a single English-centric dense decoder-only model in the $8$B class (Qwen3-8B) and to a single architectural family of attention (GQA). We do not test transferability to other dense 8B architectures, mixture-of-experts models (e.g., Llama~4 Scout, Mixtral), or different parameter scales; the findings are therefore scoped to Qwen3-8B specifically, and the generalizability of the forgetting topology to other architectures is an open question addressed in future work. We also restrict evaluation to English; multilingual forgetting may exhibit different topology.

Our causal tracing uses \textit{single-component} ablation and therefore under-measures effects that emerge from multi-component interactions. Forgetting that is distributed across a small circuit (e.g., paired induction heads, or an attention-head / MLP chain) is attributed only partially to each component in isolation. Pairwise ablation is computationally quadratic in the number of components and is out of scope for this submission; we discuss it as follow-up work.

We operate at projection granularity rather than attention-head granularity. LoRA's implementation targets full projections, so the \tglora{} target set is naturally projection-level; extending both the trace and the adapter to per-head granularity requires a custom LoRA wrapper and is left to future work.

Within the experimental matrix (1 model $\times$ 9 methods $\times$ 3 seeds $\times$ 4 stages, plus reversed-order ablation) we do not perform an exhaustive hyperparameter search for \tglora{}'s threshold $\tau$; we report a three-point sensitivity sweep (\S\ref{sec:results:ablations}) rather than an optimal-$\tau$ curve.

Finally, the reversed-order ablation tests only one alternative task permutation, not the full $4! = 24$ permutation set. We report whether \fvs{} is stable under the single most-informative ordering swap (fully reversed) rather than across all orderings.
